\MakeAbstract{
    传统音乐推荐系统多以历史播放偏好为核心信号，往往难以感知用户「此刻」的情绪状态，导致用户在情绪起伏时仍需付出较高的找歌成本。本课程设计以「面向情绪感知的本地音乐电台」为目标，设计并实现了一个名为 Mood Radio 的人机交互系统，将用户当前情绪作为推荐管线的一等输入，结合本地音乐库的情绪坐标完成实时推荐。

    系统以 Russell 提出的 Valence--Arousal 二维情绪模型为公共语义空间，统一接入摄像头表情识别、手动情绪选择、自然语言聊天、时段电台与头部手势五种输入通道；推荐管线在 SQL 层完成基于情绪距离、用户反馈、新鲜度与最近播放时间的复合打分。视觉层提供 7 套情绪主题、Aurora 流动背景与黑胶旋转封面，使整个界面随当前情绪柔和变化，强化"系统正在理解我"的反馈感。系统进一步引入匿名用户标识与心跳聚合，实现了一个隐私友好的"群体此刻"模块。

    工程实现采用 C++17 + cpp-httplib 后端、React + Vite + TailwindCSS 前端与 SQLite 持久层，并通过 onnxruntime + librosa 离线脚本完成歌曲情绪标注。课堂模拟评估表明，用户对本地隐私设计认可度较高，推荐逻辑可解释性良好，匹配/安抚双策略受到普遍欢迎。本项目较为完整地体现了多模态输入、情境感知、用户控制、可解释性与隐私保护等人机交互核心原则。
}{情绪感知，多模态交互，音乐推荐，可解释性，隐私保护}

\MakeAbstractEng{
    Most music recommenders rely on long-term listening history and rarely reason about the user's \emph{current} emotional state, leaving users with significant search cost when their mood shifts. This course project designs and implements \emph{Mood Radio}, an emotion-aware local music radio that treats the user's current mood as a first-class input and serves songs from the user's own library in real time.

    The system adopts Russell's Valence--Arousal model as the shared semantic space. Five input channels --- face, manual selection, chat, time-of-day and head gesture --- are unified into one target coordinate, and recommendation runs composite SQL scoring over emotion distance, feedback, novelty and recency. Seven mood themes, an aurora backdrop and a spinning vinyl cover morph with the current mood, while an anonymous identity layer and a heartbeat aggregator provide a privacy-respecting collective-presence view.

    The implementation pairs a C++17 / cpp-httplib backend with a React + Vite + Tailwind frontend, SQLite storage, and an onnxruntime + librosa script for offline mood labelling. A small classroom study shows that the local-first privacy design and the explainable scoring are well received. The project illustrates how multimodal input, context-awareness, explainability and privacy can be combined in one end-user HCI artefact.
}{Emotion-aware computing, Multimodal interaction, Music recommendation, Explainability, Privacy}
